
\documentclass[11pt,a4paper]{article}

\usepackage[margin=2.45cm]{geometry}
\usepackage[T1]{fontenc}
\usepackage{tgheros}
\usepackage{xcolor}
\usepackage{microtype}
\usepackage{setspace}
\usepackage{parskip}
\usepackage{fancyhdr}
\usepackage{titlesec}
\usepackage{lastpage}
\usepackage[hidelinks]{hyperref}

\renewcommand{\familydefault}{\sfdefault}
\definecolor{headinggrey}{gray}{0.22}
\definecolor{rulegrey}{gray}{0.72}
\definecolor{timestampblue}{HTML}{2F5D7C}
\definecolor{timestampback}{gray}{0.94}

\setstretch{1.08}
\setlength{\headheight}{23.2pt}
\setlength{\parindent}{0pt}
\setlength{\parskip}{0.75em}
\emergencystretch=2em

\pagestyle{fancy}
\fancyhf{}
\lhead{Lecture Transcript}
\rhead{Least Squares Regression}
\cfoot{\thepage\ of \pageref{LastPage}}
\renewcommand{\headrulewidth}{0.4pt}
\renewcommand{\headrule}{\hbox to\headwidth{\color{rulegrey}\leaders\hrule height \headrulewidth\hfill}}

\titleformat{\section}{\Large\bfseries\color{headinggrey}}{}{0pt}{}
\titlespacing*{\section}{0pt}{1.2em}{0.6em}

\newcommand{\timestamp}[1]{%
  \par\vspace{0.35em}%
  \noindent\colorbox{timestampback}{\textcolor{timestampblue}{\small\texttt{#1}}}%
  \par\vspace{0.2em}%
}

\begin{document}

\begin{center}
    {\LARGE\bfseries Lecture Transcript}\par
    \vspace{0.35em}
    {\large Least Squares Regression and Optimization}\par
    \vspace{0.6em}
    {\normalsize Speaker 1}\par
\end{center}

\vspace{0.5em}
\hrule
\vspace{1.2em}

\section*{Timestamped Transcript}

\textit{Timestamps were added from the visible Otter timestamp screenshots. Text between two timestamps is the transcript segment beginning at that timestamp.}

\timestamp{00:14}

Great pleasure to welcome Professor Roy Smith from ETH Zurich. He's come all the way here to teach you, so I think it's a great privilege, and we're all looking forward very much to learning from your lectures. I think I also posted to the course page about the workshop, which you're also welcome to attend. Keen undergraduates are part of the scope. So I think that's about it. I'll hand it over to Professor Roy. Thank you very much. Thanks very much, Jeremy, and yeah, thanks for hosting me. So I'm coming for a month as part of the visiting Erskine fellow, which you probably may or may not know about, but I'll be teaching it around for this month. You can probably tell from my accent that I'm not really Swiss, and actually, I got my undergraduate education right here in the late 1970s things



\timestamp{00:57}

have changed a little bit, but you know those lecture halls on the other side of the lower level. That's where I was educated, and e6 and e7



\timestamp{01:05}

so I've been through in exactly the same position you are. So you might be curious about how am I end up in Zurich. I was 20 years a professor in the United States, UC Santa Barbara as well.



\timestamp{01:17}

And I've worked in industry, and I've worked for NASA in this, in a variety of projects. So if you want any advice on career or options or things like this, or apart from a lot of good luck, how did that work out? Then please feel free to send me an email. I'm happy to chat and give you advice on that.



\timestamp{01:34}

So you're starting to write at the same trajectory point that I was in, so you can see how it might end up one way, for better or for worse. Anyway. Actually, I officially retired last year from ETH, because they have a mandatory retirement age, but I can't quite give it up. So this is part of the reason I'm back here. And I'm going to teach you a bit about some optimization topics. I put together some slides from a variety of courses that I've taught in the past.



\timestamp{02:01}

So my usual style is to really encourage some interaction. If you've got any questions at all, just feel free to interrupt me. This is Swiss. Tend to be a little more polite. I've also taught in Ghana and Africa these topics as well, and and they're much more excited and energetic down there, because to have a professor come from Switzerland is not all that common. So and they realize that your time is valuable, that you'll be gone in a couple of weeks. And so they ask everything they can possibly think of in the class. And so we spend the whole day.



\timestamp{02:31}

So somewhere between those two would be pretty good. So anyway, let me start. And so I wanted to start here. We have 15 minutes, I guess, talking about least squares regression, which is quite a basic thing, but I thought I would, I would introduce it because it's also a very good way of highlighting some aspects of optimization which appear not to first be least squares, but least squares or the ideas behind are buried inside and also on the other side. You can take some of the ideas you do see in big optimization, large scale optimization problems, and then look at them, realize that they're much easier to understand in terms of these squares problems where the basic things are linear, and so you can see what, what these things are doing for you. So least squares will probably talk about the next two days plus or minus, depending on how far we get. So anyway, let's let's begin. So I wanted to talk about what was the hot science topic of 1800 so we're going back to the beginning here. So, and that was to try and find an undiscovered planet in there. And about that time, well before 1800 this was the current list of planets I've left. Look at that list. I've left Earth off.



\timestamp{03:36}

So, but it was in there, and they were aware of that.



\timestamp{03:39}

So science at the time had discovered everything out to Uranus, basically. And actually, I think this is 1778 so this was not much more recently. And there why a missing planet? Well, there's something called the Titius--Bode law, which is this formula, d is the distance from the Sun to the planet in normalized units.



\timestamp{03:59}

It follows this sequence. So D, point four plus point three times two to the n, and n is minus infinity, zero, et cetera, in there.



\timestamp{04:09}

And that does give you a pretty good approximation, at least for a circular approximation, to the orbits the distance of all the planets from the sun. Now there's Earth is one. If you put n equals one there, you get Earth in there.



\timestamp{04:21}

Venus is zero and Mercury is minus infinity. But they made up this nice little sequence. And one very disturbing thing about the sequence is there's no number three in it.



\timestamp{04:29}

And so how comes no number three? There must be a planet we haven't discovered in there. And they had a pretty idea of the distance, but it doesn't tell you exactly where it is in the sky. You guess it might be in the ecliptic plane, but you wanted to find it. So the hot topic of 1800 was, well, let's try and find this planet.



\timestamp{04:46}

There isn't a planet at three. That's roughly where the asteroid belt is. There might once have been a planet at three, but it's been broken into pieces. And so the largest thing there is Ceres, which is only bit under, or just a bit over, 900 kilometers in diameter. So it's pretty small, but it's enough to see. It was very hard to see with telescope technology of 1800 but it was definitely the problem to try and find that. So that problem, I have a little history here. So in 1801 Giuseppe Piazzi, he did find Ceres, actually, and he made 24 observations of it. So got his telescope out, and an observation consists of two angles, basically. So if you sort of head to the vernal equinox out of one side, you've got a sort of a, what's called a right ascension angle to the planet. And then there's an elevation angle up. And then you see where the planet is. You record that. Record those. So each observation is two numbers, two angles, basically. And he made 24 of them over what amounted to 44 days. Weather wasn't always great. He's in Italy, and he saw, basically, 22 observations, 22 measurement points. Each one is two numbers. Now in February, Ceres tracked behind the Sun. In this because Ceres is going around one way the Earth's going around another, same direction. It wasn't going to be visible for quite some months afterwards. Now, it was so hard to find they wanted to predict when would it come out from behind the sun so they could be sure to see it. So the observations he had corresponded about nine degrees of arc in there. So that's not very much to make a prediction.



\timestamp{06:26}

So this was published by a sort of scientist and astronomical magazine editor called Zach --- Franz Xaver von Zach, editor of the \textit{Monatliche Correspondenz zur Bef\"orderung der Erd- und Himmels-Kunde} (Monthly Correspondence for the Promotion of Earth and Sky Science). He published the observational data, and all of the astronomers of the day tried to estimate when Ceres would appear on the other side of the sun, and where exactly it should look for it. So Gauss picture here, he estimated the orbit of Ceres, and Zach published the ephemeris, which is essentially the parameters that characterize the orbit, along with those of many of the other hot astronomers of the day. But Gauss was significantly different to most of the other astronomers.



\timestamp{07:01}

Then on 31st of December, so you see, this is 10 months later, and there, Zach himself got the telescope out and found Ceres within half a degree of where Gauss said it would be, but Gauss would not say how he'd done the calculation, and this, he was accused of sorcery, and it's a great way of destroying the reputation of your academic opponents. Anyway, he found it the same thing repeated in 1802 Pallas, which I think is the second largest of the asteroids, was was estimated a few points or measured, and Gauss estimated that orbit. And then he still wouldn't say an 1805 Legendre published a paper sort of on the method of least squares, but was aware that this is probably what Gauss was using, even though he wasn't saying it, and Gauss had actually been using it even ten years earlier for making comets, predictions of the trajectories of comets.



\timestamp{07:53}

So he worked on some of these techniques. So on the basis of being able to do that, he was became director of Gottingen observatory, another astronomer discovered Vesta, which is the third largest of the asteroids, and Gauss took only 10 hours to calculate its orbit. Now 10 hours to calculate the orbit from a number of X points, no calculator, no computer. It was basically actually Quill and ink on a piece of paper. And you can read his notebooks. I've had a look at some of his notebooks. You can look at them online, and there's quite a few pages, and he very rarely made any mistakes and this so, I mean, I don't think I could code up software to do that in 10 hours, let alone actually run through every calculation. And he did his calculations to about eight decimal places, and this by hand.



\timestamp{08:40}

So incredibly impressive. Actually, I should, you know, actually, I should go back here. I'm just when I'm thinking about Gauss. He came from a sort of lower middle class background. His father worked for a landowner. I think it was, and it was quite a generous person who offered to fund Gauss to go through school.



\timestamp{08:59}

And it was in school once, and the teacher was a bit bored, and it asked he wanted if the class something to occupy them for such period of time. So he asked them to sum up the numbers from one to 100 and



\timestamp{09:14}

this took about two seconds to do that.



\timestamp{09:19}

How did he do it? He



\timestamp{09:22}

didn't start go, one plus two is three. Plus three is six. He did not do that. So you said, Okay, round about 50 numbers, 49 plus 51 is 148 plus 52 is 100 so you've got 49 times 100 plus the 50 plus the 100 done. And the teacher saw, you know how he'd done this, and thought this, this kid who was quite young, is special. So paid for the education higher and higher for him. So anyway, he got to become the director of Gottingen observatory. And it was 1809, that he finally published a paper that outlined the methods of least squares in there.



\timestamp{09:59}

So actually, let me say I offered to give you some career advice.



\timestamp{10:03}

Don't do this. Don't come up with the best method the planet has seen in for hundreds of years and then take 10 years to publish it in this if you've done if you've done something pretty good, you should probably publish it fairly quickly. So anyway, what did Gauss do?



\timestamp{10:20}

All right, and I'm just out with a simple example, and we'll work up to what he actually did today. And that's fitting data to align. And this is something I remember seeing in high school. Probably seen a similar thing, and this certainly in your first year, you've got a series of points on x y plane. And I've got them, you know, x1 y1 up to x n, y, n. So I've got these, and I would like to have a line go through all of these points. So this is a, x plus b. So the things I need to figure out are, what's a and what's b. These are the sort of unknown parameters in this. Now, ideally, you've got the case, well, there's this relationship which holds. And I would call that roughly. I would call that sort of a model of what you expect in reality, particularly if these are measurements of things, this never really happens that you've got some errors in there. And typically for here, we can assume the errors are in y. So we know the x's, but Y could have some errors in them. So the real model of our system, including the errors, would add, say, an epsilon i to each one of those in there. And so we have a matrix equation vector of the y's, A's and B's. And if you look at each row of this, this is y1, is x1, times a plus b plus epsilon one. It's exactly that equation. So all of the rows are just written down in matrix form in there.



\timestamp{11:37}

And I've given some symbols for them.



\timestamp{11:40}

Y is pretty obvious. That's this one here. This matrix I'll call phi, and it actually has another name. And what we're going to do, it's called the regressor matrix.



\timestamp{11:52}

So that's v theta I'm going to use as a variable, which is my unknown parameters. So theta, unknown parameters, and my job is to find those parameters. And epsilon is just going to be a vector of errors or noises, if you like. So that's linearly. Squares is to essentially solve the problem of finding the A and B that satisfy this. Now, if you have a look and how many unknowns you've got here, you've actually, obviously a and b is two unknowns. But if you've got all of these epsilons, there's another n of those. So you've got n plus two, and you've got n equations, n simultaneous equations, if you like. So clearly, this is not unique. There are many infinite number of choices that will satisfy this. So you're looking for what is, quote, The best one. And the best one we're going to say is this one with the smallest noise in here. There's a couple of good reasons why we're going for smallest in there, and we'll get to these by the end of today. So we look for the smallest one.



\timestamp{12:53}

And we typically, as you might expect, actually look for the smallest in terms of the square of the entries here, the sum of the square, or the two norm of the entries of the epsilon, that's going to be our idea of the best possible fit, which is where the name comes from, least squares, as you might expect. It's very, very honest about that. Okay, so here's our linear fitting.



\timestamp{13:20}

We have a model equation. This is why I call it the model. This is in the previous example. This was the fact that a line goes through these points. It may be true, but that's what we assume, and that's why I say it's our model of what's going to happen in general. Let's say we've got n points, and let's say we've got D parameters to get D was two in the previous example. So that's a D dimensional vector of theta, and the regressor matrix is, of course, n by d.



\timestamp{13:48}

Now if n is bigger than D, so you've got many, many more rows, then you have an over determined case, in the sense that this equation alone is going to have many, many ways of coming up with the theta if it's exact. Now, of course, in practice, it's never exact, in that it's quite possible that n would be smaller than the number of parameters. Now, actually, it's becoming increasingly possible these days. So this is the case we normally think of it. You have a lot of measurements. You're trying to fit a line through it, or maybe a parabola or a circle. But parabola or a circle, but you have some simple form, you're a relatively small number of parameters, and you're trying to find, we've got lots of measurements of this. This case, though, comes up frequently in machine learning, and that's the case where you've got models which have many more parameters than available data. And that's become a very common thing in this and so there it's under determined there are actually an infinite number of thetas which will do the job.



\timestamp{14:48}

A lot of big question with machine learning, what's the meaning of the fact that there's an infinite number of them, and how do you handle that? So we'll get to that. But the under determined case comes up more often than you might expect, particularly in large scale data cases or large scale models. And if you think about large language models or vector support machines or essentially any of these machine learning methods, often they have a huge number of parameters, millions of parameters in language models, and not so many data measurements. Okay, but a lot of what I'll focus on here, and in a lot of engineering applications, this is the case that it's over to him, you do get more measurements than you look at parameters that you're looking for.



\timestamp{15:30}

So because this is an engineering class, I want to put this in an engineering or non engineering optimization context. One way of sorting this out is to, say, introduce these errors, epsilon. Think of the objective to be minimized as the square, the two norm squared, if you like, the sum of the squares of the elements in there. It doesn't matter so much that it's squared. It makes the mathematics easier, particularly when you come to differentiate it in that you get the same answer. But we typically write it as squared, and subject to this constraint. So your measurements y are explained by phi times theta, plus some noise, and you're trying to make that noise or error as small as possible. This is also known as an equation error form,



\timestamp{16:17}

because essentially, you would like that epsilon to be zero. That's the constraint you would like to satisfy. It's generally not possible for just picking arbitrary y. This will not be a consistent equation equal to zero in general.



\timestamp{16:31}

So we call an equation error, and this will come up a little bit later. Is sometimes the equation error corresponds to the physical intuition we have about the system. So if y was really subject to measurements X are all in here, measurement errors and X all in here are not then this is also a reasonable model of measurement noise. And so you think, OK, small measurement noise. What's the smallest measurement noise? That describes my model structure, and making it small is likely because, for measurement noise, we typically assume that I'll draw the shape. It has some probability density, density function like this, and it's zero mean, typically. And so getting it as small as possible means it's the most likely thing to have explained the result, because it's close to the middle here, and the middle is where the density function is highest. So that's motivation for making it small in terms of the probability density function. Now, if that's your entire problem description, the easiest thing to do is to take that epsilon, substitute it in here, and you end up with this optimization problem. And this one's unconstrained, makes



\timestamp{17:39}

it a little bit easier to solve.



\timestamp{17:45}

You don't have to work out KKT conditions and things like this, which I think you probably suffered with last week or so. You just, how would you solve this? You just differentiate it, set it equal to zero. Would be the obvious way to do that. So it's a little bit easier. You might not be able to do that because you may have some additional constraints that come up, you might want to constrain the theta to be in some other set. Maybe they have to be positive, or maybe there's some other convex constraint that you might want to add in, which case you can't quite simplify it like that.



\timestamp{18:13}

All right, so we'll look at this particular problem.



\timestamp{18:17}

Let's make the problem a little bit harder. We'll go through a couple of steps, and now, instead of fitting a line, I'm going to fit a linear function parameterization.



\timestamp{18:26}

Now what I mean here is that I've got my Y equals FX, but FX is not here, a, x plus b. It's a little more general. It's a linear combination of functions, H, I, of x, so and each one is multiplied by theta i there, so the H, I, you can think of them as basis functions. They tell me the basis I'm looking for. When I go back to the simple example. This one, the basis functions are really just x and one very, very simple, but there could be more general things.



\timestamp{19:04}

A common one might be polynomials. And I've given an example here where the easiest way of thinking about it is my polynomials are powers of x. So if I'm after a second order a quadratic function to fit through my data points, I would choose 1x x squared, and then find the values of theta which would multiply each of these to give me a general third order or second order polynomial. This is certainly works. In definition, it's not perhaps the best thing to do a better one, is to use Chebyshev functions. So again, the first one's 1x and the next one's 2x squared minus one. And you keep going up, you can calculate what the next one is it's 2x times the previous one minus the one before that. And I've plotted both of these options over here, and you see the red dashed lines, the regular polynomials. I think I can see 01234,



\timestamp{19:57}

so basically up to fourth order polynomial shown there. And here I'm showing also for the Chebyshev functions. And you see quite nicely that they look much more, quote, interesting, if you like, through most of this graph, they have different functions that can fill everything in the graph. If you look at zero, and you think of the powers there, the only thing that can fit data near zero is the constant line when x is close to zero. Whereas in the Chebyshev you've got multiple options. You've got this one, you've got this one, higher order ones. So the number of points, or the fit that you can get, say, across a minus one to one normalized interval, is much, much better with the Chebyshev functions, what this comes out to mean when you apply it is that the matrix fee is better conditioned than this, whereas this one might be poorly conditioned depending exactly on where the x's are. There's lots of ways of explaining sort of sort of data here or data here in this, but not much as you get down down here. So it could be you're fitting much more complicated functions than just lines, but the key point is that you're still fitting a linear combination of these functions, so you



\timestamp{21:09}

still have a linear combination to look, to look for when you do that.



\timestamp{21:14}

So we could do this, this wouldn't be particularly good for finding the orbit of a planet or a comet, or because these describe polynomials. Comets and planets don't travel along polynomial functions. They actually travel along conic sections. So we can do something similar for those.



\timestamp{21:37}

Okay, well, let's first talk about least squares. So here, I say the simplest problem. I couple of examples of what you might use for that is to find the closest fit. So this would be the smallest epsilon in here.



\timestamp{21:51}

And I said I put it in the form, or in the simplest form, it's an unconstrained optimization problem, and it has this cost function. So let me just write the cost function down, say, J. It's a function of theta, which is the variable I'm looking for, and you can see it's equal to phi theta minus y, transpose phi theta minus y. So that's the sum of the squares of the elements of that. Think of a row vector and a column vector. So you just get the sum of the squares the elements. And of course, that's actually equal to two norm v squared in there. All right, so one way of doing is to multiply this all out. And we can do this. This would be theta transpose phi, transpose phi theta minus two phi theta transpose y plus y, transpose y, y, transpose y, is not going to matter to us, but the two terms we have and here we have a quadratic in theta here and a linear term in theta buried in there.



\timestamp{22:56}

So how do we solve this? Well, this is just a function of theta. Theta might be vector valued. You can fix the theta one, theta one, theta two, theta three, if you like the components in there, we differentiate it and set it equal to zero. So differentiating that, I don't know how much experience you've had differentiating matrix or vector equations. Yes, some Okay, no problem. Yeah, it wasn't the priority. Actually, it



\timestamp{23:16}

doesn't have to be the priority. It's sort of thing that it's really nice to do once and convince yourself that it all works out in terms of where the train in terms of where the transposes go and stuff like that, and never do it again. Instead, go to this website, matrixcalculus.org, and you just gives you boxes you enter in your complicated matrix equation, and then find spirit in there. It's really great do it once, just to make sure, because there are actually a couple of different choices of notation you can use depending when you're differentiating, say, a matrix, you'd go down the columns or across the rows first. So there are a couple of different ways of doing it, but there's some more common standards, but you can find a little bit more there. So once you've done it, once convinced yourself that what I'm writing here is correct, then you can go ahead and just use that. So we differentiate this. Of course, there's a quadratic term there, so you'll expect to see, whoops. Okay, I changed color. Too bad. Two phi transpose phi. There's the term for the quadratic times theta minus two phi transpose y, and we set it equal to zero to find the minimum of this convex function in there, phi, transpose phi. Of course, the quadratic term in here, there's going to be a positive definite matrix, maybe only positive semi definite in there, if the fee is not full column rank. But we'll talk about that case a bit later on. That's what happens in machine learning. It's not, it's not positive definite. And so let me just shuffle this around a little bit and write down. Let's choose a different color here so we can write it down as phi, transpose phi. Put that in parentheses, theta equals phi, transpose y.



\timestamp{24:56}

So you see that's the linear equation that describes the solution here. And there's a couple of interesting things about that. Well, one is it's we like it so much. It's got its own name.



\timestamp{25:08}

These are called the normal equations.



\timestamp{25:11}

I don't know what the UN normal or sub normal ones are, but those are the normal ones in there. So phi transpose phi times theta equals phi, transpose y, and so the solution, the minimizers of this objective satisfy that equation. Now this equation, y minus phi, theta is not always consistent. What I mean is, for any possible y, there doesn't always exist a possible theta that fits that exactly. But this is always consistent, which is an important property, it always has a solution for any y, any solution.



\timestamp{25:47}

And so we would call this consistent. Have to be a little bit careful here, because consistent has a different meaning as well. But the one is they're looking for here is, you give me any y, and there will be a corresponding figure.



\timestamp{25:59}

Now, in the case where you've got an over determined thing. What you can always do, of course, is this is a square matrix here. It's positive definite, so it has an inverse. So you could write your optimal theta. That's called theta star is equal to phi, transpose phi, inverse phi, transpose y. And that's the solution. You have a matrix, closed form matrix form solution for your linear least squares. And you could code this up in MATLAB, do it exactly like that.



\timestamp{26:30}

I wouldn't. And that in the sense that theoretically, it's the solution, it's good for proving theorems, it's good for teaching class. It's not always good for calculating. And the reason is, you've basically squared things here, so small numbers will get smaller. Big numbers will get bigger. Then you've inverted them, those small numbers might get very big, and then you multiply them out again, so numerically, you can end up with problems in here. And in fact, if you're in the machine learning case, if you like, if you're in the under determined case, this isn't really invertible. We'll talk about that maybe tomorrow, about what you can do about that case. But this is the theoretical solution, and with a relatively small number of things, you could even code that up in MATLAB. It would be fine in there, unless you've got particularly poor data. But that's the solution to least squares problem.



\timestamp{27:21}

So yeah, that was quick, but talked about a few more things about this. So I pointed out the normal equations are always consistent. And the point I want to make here is you have to ask yourself, does phi transpose, phi inverse exist?



\timestamp{27:41}

So if it doesn't, so if so phi has, I'll call it low column rank.



\timestamp{27:51}

What I mean is it's not the full column rank. So, so think of that matrix fee. You've got tall matrix, if you like. It could be that some columns are a linear combination of the others, and then we say it loses rank. So how many linear independent columns that are in there? And if you've got D linearly independent columns, and that's all of them, then this will always exist.



\timestamp{28:12}

But if it's low column rank, then actually you've got an infinite number of solutions theta.



\timestamp{28:22}

Now they all achieve the same cost, and it is the minimum cost. It's just not a unique number. But theta probably had some meaning for you. It's the



\timestamp{28:30}

coefficient of polynomials you're looking for. Now you're saying, OK, there's an infinite number of possible polynomials, not all possible ones, but an infinite number which minimize this. That may or may not be a problem for you. If you're trying to find the orbit of a planet, it's a problem in there. So it means you haven't taken enough measurements essentially.



\timestamp{28:49}

All right, we'll deal with that tomorrow. But one way of thinking about this is that if,



\timestamp{28:58}

let's say if, phi is not full column rank,



\timestamp{29:05}

then the data is, I'll call it insufficient. There's a problem with the data. You haven't collected enough of it, or you've collected it in a special way, which is hiding something from you. Suppose your system is described by a second order polynomial or a third order polynomial. And you happen to choose points which are all in a straight line. In there, you got unlucky. You choose these points and a line goes right through them. Then there's an infinite number of things that also go through them. If you consider now the higher order polynomials. Now, every one of them minimizes the fit to going right through the points, but it's not telling you what you want to know. So the data you've taken is insufficient in this case, so you might want to take more data in there.



\timestamp{29:51}

All right, so that's basically least squares. I want to give you another sort of interpretation to this, a more geometric interpretation, and you can see that this gives you some, I think, additional insight, which we can actually use to derive different algorithms in here. So the problem we're trying to solve is minimize over theta. Notice, I think of minimize as a verb in there, min, I think of as a function. So I was right, minimize in this. This is, this is a habit I picked up from a friend of mine who knows a lot more about this than I



\timestamp{30:28}

do. So here's the function.



\timestamp{30:30}

We're trying to n times d, and we're trying to minimize that function. One way of thinking about this theta times phi. Theta is a vector. Phi is this collection of columns. So what you have here is a linear, linear combination



\timestamp{30:46}

of the columns of theta, sorry,



\timestamp{30:51}

columns of phi. It's a linear combination of those. Which ones do you want? How much of this column? How much of this column, et cetera. Do you need? So what you can do is because, OK, let's plot in n dimensional space. I'm really good when n is three, but higher than that, it's a problem. So let's and actually, because I need a higher dimension for y, let's go for n equals two. Let's plot these columns. So let's say I have a vector here, and it corresponds to, let's say fee one. Let's say the first column. Let's have another vector here corresponds to fee two. Call that second column of the two in there. Now y in here lies in a higher dimensional space. It's an n dimensional space, and n is bigger than D, in the case that we're going to focus on today. Anyway, n is bigger than D, and so y lies somewhere out of this plane. So think of these two things as describing a three dimensional or two dimensional plane here. So y is somewhere up here, out of that plane. Now I want to find the closest linear combination of these two vectors that comes as close as possible to y.



\timestamp{32:00}

What would you think it would be projection? It's exactly the projection theorem. And so, yes, you project that y down onto the plane, so it comes to that particular point here. And then from this, this tells you what the theta one vector would be. This one here. Oops, sorry, that's theta one. This one's theta, two in that and epsilon here is this. It's the error you have.



\timestamp{32:50}

So you can see that epsilon equals Y minus phi times theta, and phi times theta is this point here, phi times theta lies in this particular plane. So that is one way of thinking. What you're doing with linear least squares is projecting onto the columns of the regressor matrix, and you get to this point. So you can



\timestamp{32:55}

think of that from a separate point of view, and think about, what's the what's the property of the projection, how does it I drew this line straight down right. What was I trying to show you here?



\timestamp{33:10}

This hits at a right angle right. The error is at right angles to the plane. If it wasn't at right angles, then there would be another point. You could drop down and get even smaller error in that. So there's a right angle corner in here somewhere that's right angles to everything on the plane. So there's orthogonality. So epsilon is actually orthogonal to phi times. Let's call this theta star. That's going to be my optimal one.



\timestamp{33:38}

Yeah, this perpendicular symbol means that the two vectors are orthogonal.



\timestamp{33:43}

Okay, so let's say, let's take that and say that epsilon is orthogonal to phi times theta star. Theta star is our optimal choice, our best choice.



\timestamp{33:56}

All right, so one way of writing that these are both vectors in Rn is, if you move it



\timestamp{34:02}

perhaps, product Exactly. So let me write phi theta star. Inner product with epsilon equals zero means they're orthogonal. Are you familiar with inner products? OK, cool. That's good. So epsilon, though, phi theta star, epsilon,



\timestamp{34:19}

remember, was actually y minus phi times theta here. Actually this is for our theta star is equal to zero. And you remember the inner products are linear, phi, theta star, y minus phi, theta star, phi, theta star equals zero. Or I could start writing these things out. The inner product recall is phi theta star transpose, y minus phi theta star transpose, phi theta star zero. And if you have a look here, you can see that what I've got is theta star transpose, phi, transpose Y minus take the theta star out here. Oops, no, I didn't do it that



\timestamp{35:09}

way. So this is take too many steps. I'll get it wrong.



\timestamp{35:15}

Y, and this term is theta star transpose, phi, transpose phi theta star equals zero in



\timestamp{35:27}

there. And so the only way this can be equal to zero for theta star not equal to zero, is if phi transpose phi theta star. So this part is equal to this part.



\timestamp{35:45}

So these have to be equal, which means phi, transpose theta, theta star equals phi transpose y. And we've just seen that, right? That's the normal equations again. So the norm you get directly from the projection theorem, you get the normal equations.



\timestamp{36:02}

So that's another way of thinking of the same concept. The reason that might be useful is that, rather than thinking of expressing your optimization in terms of a cost function to differentiate and make it equal to zero, you can think about, oh, I have to find an error vector which is orthogonal to my regressor. And there are methods which derive solely from that point of view. These are called instrumental variable methods in there, and they essentially exploit the property that it's going to be orthogonal. Now they both come to the same thing, but it lets you do a different class of algorithms for some particular cases.



\timestamp{36:36}

So it's what it's worth having a look at.



\timestamp{36:40}

OK, here's some interesting plots. In here. I've given you four data sets and plotted them all versus x, y. There's some interesting things about these data sets.



\timestamp{36:54}

The first The first one is that they all have the same mean value of x. It actually happens to be Y, and I've plotted



\timestamp{37:01}

it so the mean of them, so average all of the x components, and you get nine, average all of the y components, and you get 7.5



\timestamp{37:12}

the variance in every case of x is 11.



\timestamp{37:17}

The variance of y is what, 4.125 no magic number there. And your estimate for fitting a line, a as the slope is in every case, 0.500 and b is 3.00



\timestamp{37:32}

every one of these data sets gives you exactly the same fit to a line. Now you ask yourself, should I actually really be fitting a line here? If you look at the statistical data for that, you've got the data set from a statistical point for don't look identical in there. The reason why this should be a little bit of a warning is some people just use the statistics of a data set. Is this a sufficiently rich data set? Do I have a high enough variance? Lots of ways. Okay, this looks pretty good. When you actually come and plot them, it does not look very good. Well, you look at this one and you say, OK, well, yeah, a line could be a reasonable thing here. There's quite a bit of noise, but a line might be the right thing to do. In this case, you look at this one, there's a very precise relationship between X and Y. There's just not a line. You're doing the wrong thing. You'd be much better off fitting a polynomial to that. In there, this one, there is a line, but one outlier messed it up for you and made it appear incorrect.



\timestamp{38:30}

Now, this one, there's no relationship between y and x, but again, one outlier gave you the impression there is and all of these, you get the same line in this but you should not be using it for these three problems in there. You'd probably use it for this one. So this is an example published by Anscombe in 1973 and it was called graphs in statistical analysis. The point he was making was have a look at your data visually before you decide that you're going to fit a least squares thing to it, and see what it actually does. Unfortunately, that message, I think, is kind of lost these days. People just put it into some big learning algorithm. They come up with a model, and they've no idea whether they've done the right thing in this. Three of these cases, you've done exactly the wrong thing. One of them is probably reasonable up in here. So think about your data in this. And now, when you've got a lot of it, it's really hard to do, really hard to visualize in this. But in 1973 it was clear to Anscombe that people were misapplying a lot of these methods for fitting data.



\timestamp{39:30}

Ok? And the reason it becomes relevant is that these days, a lot of optimization problems are driven by data in this. It used to be that maybe 1020, 30 years ago and more, until the 1960s that a lot of the optimization problems were usually for a design or scheduling a process, where it wasn't experimental data you were typically looking at. It was some features of a process, a model of a system, and this, or maybe it was scheduling for your production plant. Is my oil going to arrive through the Straits of Hormuz to fill up the now defunct plant at Marsden point. I know someone who actually she spent most of it. She was a student here at St Thomas me and most of her engineering career was spent solving the linear program in order to schedule the optimal fracking. And well, not fracking, but the refining of the different oil products. That would be a tricky, tricky problem to be dealing with right now, and as that's been shut down, there's no need. You just have to pay more. Okay. So here we have, okay. Let me get back to the original problem.



\timestamp{40:35}

So here, the point is that I would say that more and more in the last decade, optimization problems are actually being applied to data that's come from measure. Come from measurements. This could be movie recommendations. It's one of the big ones to begin with, and it's but it's data that was collected from a population, and quite literally, usually, a population of people, and this about certain behaviors and habits, and try and fit models to this data. So as such, the data now contains noise in this and so this is why we're looking at optimization methods that are designed to be handled handled noisy data or stochastic data.



\timestamp{41:08}

Okay, let's go back to Gauss again and see if we can interpret some of this in terms of what he did.



\timestamp{41:15}

So planets essentially are conic sections in that you can go to Wikipedia, there's lots of good descriptions of conic sections, but conic sections. But basically, if you think of having a cone like this, you can slice it right through. Here you get a circle. You can slice it on an angle, you'll get an ellipse.



\timestamp{41:32}

Basically, you slice it through here you'll get a parabola. Oops, no, you got a hyperbola that one.



\timestamp{41:39}

So you'll get an orbit that comes through here. You can slice it also. This is a terrible drawing. Slice it vertically through there, and you'll get hyperbole on this, but they will generate basically the curves that describe the movement some comets only go through the solar system once they're actually traversing. Hyperbole. The ellipses are the ones that are in orbit in this so you can parameterize all of these conic sections by an equation like this. It's a linear function, or linear parameters times what are called monomials of second order. So x squared x times y, y squared x, y and a constant one in there, at least one of A, B or C has to be not equal to be not equal to zero. Otherwise you get something degenerate, a line, or a, worst case, just a point in there, if you went right through the origin. So these are the various things you go through. The parabola doesn't have a constant. It's the one that's right at the boundary of whether it comes back in or not. For an ellipse, OK. So you could actually set up a linear least squares problem. Listen, find A, B, C, D, F, actually, there are only five of them, because you notice, you could scale by any constant and it would still be satisfied. So there's only five independent parameters there. You could set f to be one. It wouldn't get you the parabola, but, oh well, we wouldn't do too badly, all right. But the problem with this is that I could write down the equation error. I could set up at least squares by, say, making that equal to epsilon. But that epsilon doesn't actually have a physical meaning. It's just the equation error. How far out is this equation? Now, because the x and y come from measurements, particularly they come from Piazza, pointing this telescope at the sky, and the angle I'm pointing this telescope is such and such, crude telescope called measurement of angle, the noise is actually in the basis functions in there, and that causes a problem in there. So I've got two or three minutes, so maybe we'll continue some of this a little later. But you could write it down as an equation like this, and you could solve that. But if you've got noise in the x and the y, if you really write out, okay, here's how you might find the quote best conic section that isn't a parabola in here, scale by the F. And then you write down theta, 132, 3334 35 et cetera. And then you set up with the nearly squares problem in this. And then this would be your regressor with the columns of monomials that come from your data. Theta are the five parameters. One is the constant in here, and here's epsilon. You minimize it standard linearly, squares in there. You can solve it exactly this way I've showed you so far. But in practice, what's really the case is there's some error in epsilon one, and you square it. Now, even if this is a random error, once you squared it, you've got.



\end{document}
